% MAIN SNAPSHOT — LSPV/RMV::SYNTHESIS-0025::2026-10-12T15:47:30Z
\documentclass[11pt]{article}
\usepackage{fontspec}
\usepackage{geometry}
\usepackage{amsmath,amssymb,mathtools}
\usepackage{booktabs,longtable,array,tabularx}
\usepackage{enumitem}
\usepackage[hidelinks]{hyperref}
\geometry{margin=1in}
\setlist{nosep}
\newcommand{\OPEN}{\textsf{OPEN}}
\newcommand{\EXACT}{\textsf{PROVED/EXACT}}
\newcommand{\KNOWN}{\textsf{KNOWN THEOREM}}
\newcommand{\COMP}{\textsf{COMPUTATIONAL}}
\newcommand{\RET}{\textsf{FALSIFIED/RETIRED}}
\newcommand{\CAND}{\textsf{CANDIDATE}}
\newcommand{\same}{\textnormal{\textsc{same\_realization}}}
\title{Universal-Core Synthesis: Semantic Declassification Does Not Supply the Missing Transfer}
\author{LSPV-5 / MAIN synthesis}
\date{12 October 2026}

\begin{document}
\sloppy
\maketitle

\section*{Executive status}
\textbf{P versus NP remains open.}  This run closes one proposed repair to the
trace-extraction program, not the complexity separation.  Requiring a retained
``derived'' field to satisfy strict source noninterference erases every
nonconstant compatibility bit.  Allowing that bit by declassification forces a
choice: a merely syntactic tag is forgeable by extensionally equivalent COPY,
RECOMP, INTERP, and delay wrappers, while a sound and complete semantic
validator computes the released bit itself.  That validation cost is real, but
it is not evidence that the original SAT execution exposed or paid for the
bit.

\paragraph{Strict status.}
\textsf{NONINTERFERENCE-DECLASSIFICATION-TRILEMMA1} is \EXACT{} as an
elementary theorem about deterministic views.  Its exhaustive two-source-bit
audit is \COMP{}.  The universal bridge
\textsf{TYPED-DECLASSIFICATION-AS-WEAKER-BRIDGE1} is \RET{} at this level of
generality.  Construction-specific information-flow type systems remain
possible, but no theorem transfers them to every uniform deterministic SAT
decider.  \textsf{SAME-TRACE-PAID-SEMANTIC-TRANSFER1} remains \OPEN{}.

\section*{Recovery preflight}
The recovery prompt named
\texttt{research/main/2026-10-12\_0850\_MAIN\_SYNTHESIS-0023.tex} as the
newest verified snapshot.  A fresh canonical read instead found the later
verified file
\texttt{research/main/2026-10-12\_2349\_MAIN\_SYNTHESIS-0024.tex}, including
its required header and final sentinel.  Hence the old recovery condition is
resolved and no duplicate was written.  No exact body for a still-missing
later snapshot was available; none was reconstructed.  The stale predecessor
observation is best diagnosed as state lag rather than a current repository
write failure.

\section*{Exact model, encoding, and clean core}
Let $\langle F\rangle$ be the ordinary self-delimiting encoding of a 3-CNF
$F$, and let
\[
 N=|\langle F\rangle|.
\]
For a partial assignment $\alpha$, define the compatibility relation and its
one-bit decision quotient by
\[
 H(F,\alpha)=\{\beta: \alpha\cup\beta\models F\},
 \qquad
 h(F,\alpha)={\bf 1}[H(F,\alpha)\ne\varnothing].
\]
The clean-core family used for the universal target is
\[
 \mathcal C_N=\{(F,\alpha): |\langle F,\alpha\rangle|=N,\ F\text{ is a
 3-CNF and }\alpha\text{ is a self-delimiting partial assignment}\}.
\]
This family deliberately includes Horn, affine/Tseitin, 2SAT-encoded,
bounded-treewidth, padded, and redundantly encoded instances.  Excluding those
tractable controls by definition would not prove a lower bound for arbitrary
3SAT algorithms.

Fix a uniform deterministic polynomial-time SAT decider $A$.  For input
$x=\langle F,\alpha\rangle$, let
\[
 \operatorname{Tr}_A(x)=(\mu_0,\ldots,\mu_{T_A(N)})
\]
be the exact realized mutable-state trace, including preprocessing and the
one-bit answer.  A retained view is a deterministic map
$Q_A(x)=q(\operatorname{Tr}_A(x))$.  Every extraction, validation, observation,
and reconstruction step is charged in encoded length $N$.

\paragraph{Target theorem, stated without circularity.}
The desired bridge would say that for every such $A$ deciding 3SAT, there is a
uniform polynomial-time observation rule and extractor, operating on the
\emph{same realized execution}, whose total cost can be lower-bounded in an
unconditional endpoint model strongly enough to contradict the assumed
polynomial bound.  Merely asserting that no polynomial $A$ exists restates
$\mathsf P\ne\mathsf{NP}$ and is not a premise.  No such bridge is proved here.

\section*{Candidate under attack}
The previous synthesis isolated full-trace semantic factorization as the most
plausible remaining observation repair.  The newest peer reports sharpened a
natural proposal:
\begin{quote}
Tag retained trace fields as \emph{derived} and require an information-flow
policy that prevents direct copies of source input.  Then use the derived
field as the solver-independent compatibility exposure.
\end{quote}
Call this
\textsf{TYPED-DECLASSIFICATION-AS-WEAKER-BRIDGE1}.  It appears to be weaker
than requiring all of $H$, but the exact quantifiers reveal a trilemma.

\section*{Theorem: NONINTERFERENCE-DECLASSIFICATION-TRILEMMA1}
Let $X$ be any source domain, let $h:X\to\{0,1\}$ be nonconstant, and let
$Q:X\to Y$ be a deterministic retained view.

\begin{enumerate}
\item If $Q$ satisfies strict source noninterference,
\[
  \forall x,x'\in X\quad Q(x)=Q(x'),
\]
then no decoder $D:Y\to\{0,1\}$ satisfies $D(Q(x))=h(x)$ for all $x$.
\item If a policy declassifies $h$, then a syntactic ``derived'' tag alone does
not certify provenance: an equivalent wrapper can compute a source-dependent
bit by COPY, RECOMP, INTERP, or delayed copying and attach the same tag.
\item If instead a semantic validator $V:X\times\{0,1\}\to\{0,1\}$ is sound
and complete,
\[
 V(x,b)=1 \Longleftrightarrow b=h(x),
\]
then $h$ is recovered with at most two validator calls:
\[
 h(x)=\text{the unique }b\in\{0,1\}\text{ for which }V(x,b)=1.
\]
Thus semantic validation is computationally equivalent to the released bit up
to constant-factor oracle use.  Its charged cost need not occur in, or be
recoverable from, the original execution of $A$.
\end{enumerate}

\paragraph{Proof.}
For (1), $D\circ Q$ is constant because $Q$ is constant, contradicting that
$h$ is nonconstant.  For (2), a tag is metadata rather than a semantic
property; extensionally equivalent programs can produce the same tagged bit
through distinct dataflow.  For (3), exhaust both Boolean candidates and use
soundness and completeness to select the unique accepted value.  This proves
all three statements. \hfill$\square$

\paragraph{Consequence for compatibility.}
Take $X=\mathcal C_N$ and $h(F,\alpha)={\bf 1}[F|_\alpha\text{ is
satisfiable}]$.  Whenever $h$ is nonconstant on the chosen slice, strict
noninterference destroys it.  A polynomial semantic validator would compute
that SAT bit.  Under the counterfactual $\mathsf P=\mathsf{NP}$ this is
perfectly consistent; it yields no unconditional contradiction.  Without
that counterfactual, claiming the validator must be superpolynomial is simply
the missing separation in another form.

\section*{Bounded reproducible audit}
The script
\texttt{experiments/review/noninterference\_declassification\_trilemma\_synthesis.py}
enumerates all $16$ Boolean functions on a two-bit source domain.  It performs
four exact checks:
\begin{itemize}
\item strict noninterference retains exactly the two constant functions;
\item all $14$ nonconstant functions are therefore erased;
\item the sound/complete equality validator reconstructs every function on
every source, giving $64/64$ successful checks;
\item one untrusted syntactic tag is compatible with all $16$ functions.
\end{itemize}
The harness also records $168$ hostile-control cells from six wrappers, seven
formula families, and four encodings.  Those cells are not performance data
and no asymptotic conclusion is extrapolated from them.  They make explicit
that the logical result is independent of solver base, clause order, padding,
and whether a family is tractable by implication graphs, Horn propagation,
Gaussian elimination, or bounded-treewidth dynamic programming.  The result
is stored at
\texttt{results/raw/review/2026-10-12\_noninterference\_declassification\_trilemma.json}
with canonical payload SHA-256
\texttt{5581d5656ef758cf99e37a4783a5d8ee9161f93e5e813a28d08ea962731c28db}.

\section*{Latest-peer reconciliation}
\begin{center}
\begin{tabularx}{\textwidth}{p{0.17\textwidth}Xp{0.18\textwidth}}
\toprule
Artifact & Verified contribution and synthesis consequence & Status here\\
\midrule
LSPV-1 DERIVEDLEAK1 & Input-erased mutable traces do not prevent an internally
stored COPY from leaking source dependence.  This motivates semantic rather
than merely field-level provenance. & Retained\\
LSPV-2 PRODSTREAM1 & Independent products need not retain an answer vector;
streaming can reduce peak space to the maximum component cost plus counters.
Semantic multiplicity is not a memory lower bound. & Retained\\
LSPV-3 FIELDMASK1 & A fixed field mask either retains a COPY-capable value
field or drops it and erases the derived bit.  The present theorem shows why
``semantic declassification'' does not automatically repair that trilemma. &
Strengthened\\
LSPV-4 MODELBRIDGE-0027 & Full-trace factorization may simply decode the input
from the initial configuration and recompute the target, without proving paid
exposure by the original solver. & Decisive control\\
\bottomrule
\end{tabularx}
\end{center}

The reports agree on one missing arrow: a property of a representation,
trace field, proof system, or validator does not transfer to every computation
without a theorem binding the property to the same realized execution and
charging the original algorithm for it.

\section*{Faithful comparison with actual main4 and main5}
The repository entry points were read directly.  \texttt{main4.tex} identifies
itself as the ARC achievable-region manuscript; \texttt{main5.tex} identifies
itself as the WRAPPER construction manuscript.  Both accumulate explicit
partial constructions, charged controllers, exact finite audits, and hostile
controls.  Both repeatedly preserve all-input acquisition/coverage and the
separation as OPEN.  Their conditional $K$-elimination and promise-domain
results are construction-specific.  They do not state that every uniform
polynomial SAT decider exposes $H$, uses a named trace field, retains an answer
vector, or realizes the ARC/WRAPPER representation.  This synthesis does not
silently redefine either manuscript as a universal computation lower bound.

\section*{Barrier and model audit}
\paragraph{Relativization and algebrization.}
The trilemma is representation-level and relativizes: with an oracle, a
semantic validator can query the oracle just as the original decider can.
Therefore the theorem is not a nonrelativizing separation technique.  No
arithmetization or low-degree extension is used, so no claimed escape from the
algebrization barrier exists.

\paragraph{Natural proofs.}
No large constructive circuit property is supplied.  The result neither
invokes nor evades the Razborov--Rudich barrier.  Treating ``has an exposed
compatibility bit'' as such a property would first require a transfer from
program executions to a circuit class; that transfer is absent.

\paragraph{Rice-style semantic barriers.}
Nontrivial extensional program properties are undecidable in unrestricted
models.  Our finite-time, fixed-machine setting is not a direct application of
Rice's theorem, but it reinforces the need to distinguish a decidable bounded
trace check from universal semantic provenance across all equivalent programs.
Translation validation can verify a particular compilation or execution;
without a universal compiler contract it does not constrain arbitrary SAT
deciders.

\paragraph{Nonuniformity and arbitrary polynomial exponents.}
The target quantifies over uniform deterministic machines, each with its own
constant exponent $k_A$.  No single time cap $N^k$ is imposed on all machines.
Circuit lower bounds, branching-program lower bounds, and time--space
tradeoffs remain meaningful only in their stated models; none is promoted here
to an arbitrary-computation theorem.

\paragraph{One-bit output.}
A SAT decider is required to output one bit.  It need not externalize a witness,
compatibility table, answer vector, cut transcript, or proof.  Any theorem that
requires these objects must prove a solver-to-interface transfer rather than
infer it from semantic cardinality.

\section*{P=NP counterfactual audit}
Assume 3SAT has a uniform deterministic polynomial-time decider $A$.
Self-reduction then gives polynomial-time witness recovery, and the bit
$h(F,\alpha)$ is computable in polynomial time.  A semantic declassification
validator can therefore recompute $h$ in polynomial time.  This is
model-matched and consistent with the assumption; it does not contradict a
known unconditional lower bound.  To obtain a separation one would still need
either:
\begin{enumerate}
\item an unconditional superpolynomial lower bound for the same general
uniform deterministic model, or
\item a transfer forcing $A$ itself, on the same realization, into a restricted
model for which such a lower bound is known.
\end{enumerate}
Neither arrow is presently available.  Any premise that simply forbids the
polynomial validator would smuggle $\mathsf P\ne\mathsf{NP}$ into the proof.

\section*{Proof spine and missing arrows}
The surviving spine is:
\[
\begin{array}{c}
\text{arbitrary uniform polynomial SAT decider }A\\
\Downarrow\quad\text{\OPEN: same-realization semantic transfer}\\
\text{a construction-independent charged object or restricted model}\\
\Downarrow\quad\text{\OPEN: model-matched unconditional lower bound}\\
\text{contradiction.}
\end{array}
\]
The present theorem removes the shortcut ``declare the object derived and
validate the declaration.''  Strict noninterference loses the bit; syntactic
declassification is forgeable; semantic validation recomputes the bit and must
be separately charged.  The next candidate must therefore be narrower and
construction-specific:
\begin{quote}
\textsf{CONSTRUCTION-SPECIFIC-SEMANTIC-COMPILER1}: exhibit an explicit uniform
compiler whose generated SAT deciders carry a sound semantic declassification
certificate in their same realized traces, then prove a lower bound for that
exact compiled model.
\end{quote}
Even success would constrain the compiler image, not every SAT decider, unless
an independent universality theorem is added.

\section*{Stale or overstated claims}
The following universal readings are \RET{}:
\begin{itemize}
\item every polynomial SAT decider must materialize $H(F,\alpha)$;
\item every such decider must retain a source-dependent ``derived'' trace
field distinguishable by a fixed mask or syntactic type;
\item independent semantic components force a simultaneous answer vector or
sum-type peak memory;
\item a validator external to the execution proves that the execution paid for
the validated fact;
\item a lower bound for a proof system, circuit family, branching program,
streaming interface, or named construction automatically lower-bounds arbitrary
uniform deterministic computation.
\end{itemize}

\section*{Current ledger}
\begin{center}
\begin{tabular}{p{0.58\textwidth}p{0.20\textwidth}}
\toprule
Item & Status\\
\midrule
Noninterference/declassification trilemma & \EXACT{}\\
Finite $16$-function audit & \COMP{}\\
Typed declassification as a general weaker bridge & \RET{}\\
main4/main5 partial construction claims & As stated; no universal transfer\\
Same-realization paid semantic transfer & \OPEN{}\\
Model-matched unconditional contradiction & \OPEN{}\\
P versus NP & \OPEN{}\\
\bottomrule
\end{tabular}
\end{center}

\section*{Peer handoff}
LSPV-1 should formalize the strongest provenance property invariant under
COPY, RECOMP, INTERP, delay, and state recoding without becoming constant.
LSPV-2 should attack any claimed validation lower bound with streaming
recomputation and one-bit output.  LSPV-3 should test semantic field types
against wrappers that preserve extensional behavior but rewrite dataflow.
LSPV-4 should demand that any validator receive only the same realized trace,
not the original formula, initial configuration, or an alternate run.  LSPV-5
should accept no synthesis arrow that charges an external validation procedure
to the original decider without a proved reduction in the identical model.

\section*{Primary references checked}
\begin{itemize}
\item H. G. Rice, ``Classes of recursively enumerable sets and their decision
problems,'' Transactions of the AMS 74 (1953), 358--366,
\url{https://www.ams.org/journals/tran/1953-074-02/S0002-9947-1953-0053041-6/}.
\item A. Pnueli, M. Siegel, and E. Singerman, ``Translation validation,''
TACAS 1998, LNCS 1384, 151--166,
\url{https://doi.org/10.1007/BFb0054170}.
\item A. Sabelfeld and D. Sands, ``Dimensions and principles of
declassification,'' CSFW 2005, 255--269,
\url{https://doi.org/10.1109/CSFW.2005.15}.
\item M. R. Clarkson and F. B. Schneider, ``Hyperproperties,'' Journal of
Computer Security 18(6), 2010, 1157--1210,
\url{https://doi.org/10.3233/JCS-2009-0393}.
\item S. Aaronson and A. Wigderson, ``Algebrization: a new barrier in
complexity theory,'' ACM TOCT 1(1), 2009,
\url{https://doi.org/10.1145/1490270.1490272}.
\item A. A. Razborov and S. Rudich, ``Natural proofs,'' Journal of Computer
and System Sciences 55(1), 1997, 24--35,
\url{https://doi.org/10.1006/jcss.1997.1494}.
\item T. Baker, J. Gill, and R. Solovay, ``Relativizations of the P=?NP
question,'' SIAM Journal on Computing 4(4), 1975, 431--442,
\url{https://doi.org/10.1137/0204037}.
\item L. Fortnow, R. Lipton, D. van Melkebeek, and A. Viglas,
``Time-space lower bounds for satisfiability,'' Journal of the ACM 52(6),
2005, 835--865, \url{https://doi.org/10.1145/1101821.1101822}.
\end{itemize}

% END_OF_VERIFIED_MAIN::SYNTHESIS-0025
\end{document}
